\subsection{Serverseitige PostgreSQL-Persistenz mit SQLAlchemy und Alembic}
\label{subsec:persistenz-postgresql}

PostgreSQL ist optional und kein eigenes Profil. \texttt{postgres} setzt
\texttt{database} und damit ein Backend voraus; Web-only und Desktop-local
sind ausgeschlossen. Die Datenbankfähigkeit umfasst SQLAlchemy-Engine,
Sessions und Alembic. PostgreSQL ergänzt Psycopg~3, URL-Vorgaben und
Integrationstests und bleibt eine getrennte Laufzeiteinheit
\parencite{postgresql2026docs,kleiveist2026database}.

Die Engine entsteht erst beim Zugriff und prüft Verbindungen vor der Nutzung;
Sessions werden kontrolliert geschlossen. Die leere deklarative \texttt{Base}
ist ein Erweiterungspunkt für Fachmodelle
\parencite{sqlalchemy2026engines}.

Alembic nutzt dieselbe Metadatenbasis für Online- und Offline-Migrationen. Ohne
Fachmodelle gibt es keine initiale Revision. Weder \texttt{create\_all()} noch
automatische Migrationen beim FastAPI-Start sind implementiert; Änderungen
erfordern explizite Alembic-Kommandos
\parencite{alembic2026tutorial,kleiveist2026database}. \texttt{/api/health}
bleibt datenbankunabhängig, \texttt{/api/ready} prüft bei aktivierter
Datenbankfähigkeit die Erreichbarkeit.
\beginnerinput{workspace/statement/300_architektur/370}
